\documentclass[11pt,a4paper]{article}

% ---------------------------------------------------------
% PREAMBLE (stable, warning-free, Overleaf-safe)
% ---------------------------------------------------------
\usepackage[utf8]{inputenc}
\usepackage[T1]{fontenc}
\usepackage{lmodern}
\usepackage{amsmath, amssymb, amsfonts}
\usepackage{bm}
\usepackage{geometry}
\usepackage{graphicx}
\usepackage{hyperref}
\usepackage{cite}
\usepackage{authblk}
\usepackage{physics}
\usepackage{mathtools}

\geometry{margin=2.4cm}

\title{\textbf{Companion LXVII: Cross-Covariance of Persistence Diagrams in the Relaxation-Driven Cyclic Cosmology}}
\author{Michael Lehmann \\ \texttt{mi.lehmann@gmx.de}}
\date{April 2026}

\begin{document}
\maketitle

\begin{abstract}
The cross-covariance of persistence diagrams quantifies the correlated topological 
structure between different realizations, smoothing scales, or redshift slices of a 
cosmological field. In weak-lensing analyses, this covariance captures how 
connected components, loops, and void-like structures co-vary across the 
gravitational field. In the standard cosmological model, cross-covariance reflects 
the nonlinear matter distribution and the projection of large-scale structure. In 
the Relaxation-Driven Cyclic Cosmology (RDCC), the scale-dependent evolution of the 
gravitational potential modifies the correlation structure of topological features 
in a characteristic way. This Companion develops a continuous description of 
cross-covariance in RDCC, deriving how the relaxation scale influences correlated 
topology, persistence-lifetime coupling, and the observational signatures in 
weak-lensing surveys. The resulting framework provides a distinctive probe of the 
RDCC gravitational field through correlated topological structure.
\end{abstract}
% ---------------------------------------------------------
\section{Introduction}

Persistence diagrams encode the birth and death of topological features in 
weak-lensing convergence fields. While ensemble variance measures the spread of 
topological features across realizations, the cross-covariance quantifies how 
topological structures co-vary between different realizations, smoothing scales, or 
redshift slices.

In the standard cosmological model, cross-covariance reflects the nonlinear matter 
distribution and the projection of large-scale structure. In RDCC, the 
gravitational potential evolves differently. The relaxation mechanism suppresses 
the decay of small-scale modes and enhances the coherence of large-scale modes. 
This modifies the correlation structure of topological features in a scale-dependent 
manner, producing a distinctive signature in cross-covariance analyses.
% ---------------------------------------------------------
\section{Cross-Covariance of Persistence Diagrams}

Given two ensembles of persistence diagrams $\{D_{i}\}$ and $\{E_{i}\}$ with 
Fréchet means $D^{*}$ and $E^{*}$, the cross-covariance is defined as

\begin{equation}
    \mathrm{Cov}(D,E)
    = \frac{1}{N}
      \sum_{i=1}^{N}
      W_{p}(D_{i}, D^{*})
      W_{p}(E_{i}, E^{*}),
\end{equation}

where $W_{p}$ is the $p$-Wasserstein distance.

This covariance measures the degree of correlated topological structure.
% ---------------------------------------------------------
\section{RDCC Modification of Persistence Diagrams}

In RDCC, the convergence evolves as
\begin{equation}
    \kappa_{\mathrm{RDCC}}(\ell)
    = \kappa(\ell)
      \left[
        1 - \chi_{\kappa}
        \left(\frac{\ell}{\ell_{\mathrm{rel}}}\right)^{\alpha}
      \right].
\end{equation}

This modifies the birth and death thresholds of topological features:

\begin{align}
    b_{\mathrm{RDCC}} &= b \left[ 1 - \chi_{b} (\ell/\ell_{\mathrm{rel}})^{\alpha} \right], \\
    d_{\mathrm{RDCC}} &= d \left[ 1 - \chi_{d} (\ell/\ell_{\mathrm{rel}})^{\alpha} \right].
\end{align}

The RDCC ensembles become
\begin{equation}
    \{ D_{i,\mathrm{RDCC}} \}, \quad
    \{ E_{i,\mathrm{RDCC}} \}.
\end{equation}
% ---------------------------------------------------------
\section{Cross-Covariance in RDCC}

The RDCC cross-covariance is

\begin{equation}
    \mathrm{Cov}_{\mathrm{RDCC}}(D,E)
    = \frac{1}{N}
      \sum_{i=1}^{N}
      W_{p}(D_{i,\mathrm{RDCC}}, D^{*}_{\mathrm{RDCC}})
      W_{p}(E_{i,\mathrm{RDCC}}, E^{*}_{\mathrm{RDCC}}).
\end{equation}

To first order in the relaxation parameter,

\begin{equation}
    \mathrm{Cov}_{\mathrm{RDCC}}(D,E)
    = \mathrm{Cov}(D,E)
      \left[
        1 - \chi_{\mathrm{Cov}}
        \left(\frac{\ell}{\ell_{\mathrm{rel}}}\right)^{\alpha}
      \right].
\end{equation}

This modification reflects the relaxation-driven change in correlated topology.
% ---------------------------------------------------------
\section{Correlated Topological Signatures of RDCC}

The relaxation mechanism enhances the coherence of large-scale modes. Cross-
covariance therefore exhibits:

\begin{itemize}
    \item enhanced correlation of large-scale topological features,
    \item suppressed correlation of small-scale features,
    \item a characteristic turnover near $\ell_{\mathrm{rel}}$,
    \item modified redshift evolution of topological coupling,
    \item altered cross-scale persistence-lifetime correlations.
\end{itemize}

These signatures provide a direct probe of the RDCC gravitational field through 
correlated topological structure.
% ---------------------------------------------------------
\section{Conclusion}

The cross-covariance of persistence diagrams in the Relaxation-Driven Cyclic 
Cosmology reflects the scale-dependent evolution of the gravitational potential and 
the nonlinear matter distribution. The relaxation mechanism modifies the 
correlation structure of topological features in a scale-dependent manner, 
producing distinctive signatures in weak-lensing surveys. These effects provide a 
direct observational window into the relaxation scale and the temporal structure of 
the RDCC gravitational field. The present Companion establishes the theoretical 
foundation for future studies of correlated topology, nonlinear structure, and the 
interplay between light and gravity in RDCC.

\bibliographystyle{apsrev4-2}
\nocite{*}
\bibliography{references}


\end{document}
